\section{Investment Optimization}
This section outlines the investment optimization framework used to evaluate the financial viability of deploying a Battery Energy Storage
\subsection{Objective}
The primary objective of the investment optimization is to evaluate and rank the financial viability of deploying a Battery Energy Storage System (BESS) across various European countries and technical configurations. This is achieved by performing a 10-year Discounted Cash Flow (DCF) analysis to determine the Net Present Value (NPV) and the Levelized Return on Investment (ROI) for each scenario.

\subsection{Financial Modeling Approach}
As established, the core of our analysis is a DCF model that correctly aligns nominal cash flows with a nominal discount rate. This approach ensures that the effects of inflation are accounted for consistently and accurately.

\begin{itemize}
    \item \textbf{Cash Flows:} We project future profits by growing the initial year's profit (from the operational optimization) at the country-specific inflation rate. This creates a stream of \textbf{nominal cash flows}.
    \item \textbf{Discount Rate:} We use the country-specific Weighted Average Cost of Capital (WACC) as the \textbf{nominal discount rate}.
\end{itemize}

\subsection{Step-by-Step Calculation}
The analysis follows a three-step process for the results from 45 scenarios (5 countries $\times$ 9 configurations).

\begin{table}[h]
\centering
\begin{tabular}{|l|l|l|l|}
\hline
WACC & Value & & \\
\hline
Inflation Rate & value & & \\
\hline
Discount Rate & Value & & \\
\hline
Yearly Profits (2024) & Value & & \\
\hline
\multicolumn{4}{|c|}{} \\
\hline
Year & & Initial Investment [kEUR/MWh] & Yearly profits [kEUR/MWh] \\
\hline
& 2023 & & \\
\hline
& 2024 & & \\
\hline
& 2025 & & \\
\hline
& 2026 & & \\
\hline
& 2027 & & \\
\hline
& 2028 & & \\
\hline
& 2029 & & \\
\hline
& 2030 & & \\
\hline
& 2031 & & \\
\hline
& 2032 & & \\
\hline
& 2033 & & \\
\hline
\end{tabular}
\end{table}

\subsubsection{Step 1: Projecting Nominal Profits}
First, for each country among the five, we take the BEST annual profit of this country (that is, the highest annual profit of nine different configuration settings) calculated by the Pyomo model for the year 2024, denoted as $\Pi_{2024}$, and project it over a 10-year horizon from year $y=1$ (2024) to $y=10$ (2033). The nominal profit for any future year $y$ is calculated by applying the country's annual inflation rate, $\pi$.

\begin{equation}
\Pi_y = \Pi_{2024} \cdot (1 + \pi)^{y-1}
\end{equation}

\subsubsection{Step 2: Calculating Net Present Value (NPV)}
Next, we calculate the NPV by summing the discounted values of all future nominal profits and subtracting the initial Capital Expenditure (CAPEX). The CAPEX is given as \$200 per kWh. For a 1 MW BESS with a nominal energy capacity $E_{nom}$ (in MWh), the CAPEX is $200 \times E_{nom} \times 1000$.

The WACC is denoted by $i$.

\begin{equation}
\text{NPV} = \left( \sum_{y=1}^{10} \frac{\Pi_y}{(1 + i)^y} \right) - \text{CAPEX}
\end{equation}

\subsubsection{Step 3: Calculating Levelized Return on Investment (ROI)}
Finally, we calculate the Levelized ROI as specified in the competition guidelines. This metric annualizes the return relative to the initial investment.

\begin{equation}
\text{Levelized ROI} (\%) = \frac{\text{PV}(\text{Total Profits})}{\text{CAPEX} \times \text{Lifetime}} \times 100
\end{equation}
Where:
\begin{itemize}
    \item $\text{PV}(\text{Total Profits})$ is the first term in the NPV equation: $\sum_{y=1}^{10} \frac{\Pi_y}{(1 + i)^y}$.
    \item $\text{CAPEX}$ is the initial investment cost.
    \item $\text{Lifetime}$ is the project duration, which is 10 years.
\end{itemize}

\subsection{Example Calculation}
Let's consider a hypothetical scenario for a 1 MW BESS:
\begin{itemize}
    \item \textbf{BESS Configuration:} C-rate = 0.5 C $\implies E_{nom} = 2$ MWh
    \item \textbf{Initial CAPEX:} $200 \, \text{EUR/kWh} \times 2000 \, \text{kWh} = 400,000 \, \text{EUR}$
    \item \textbf{Country Parameters:} WACC ($i$) = 7.0\%, Inflation ($\pi$) = 2.0\%
    \item \textbf{Simulated Profit (2024):} $\Pi_{2024} = 50,000 \, \text{EUR}$
\end{itemize}


The Levelized ROI is then:
\[
\text{Levelized ROI} = \frac{395,437.12}{400,000 \times 10} \times 100 = \frac{395,437.12}{4,000,000} \times 100 = 9.89\%
\]
This result indicates that, for this specific scenario, the project yields an average annual return of 9.89\% on its initial investment in present value terms, but it does not break even (NPV < 0).
